\section{Stage 002: Matter-Stress Inter-Defect Force Assembly}

\stagefield{Part / anchor}{Part II -- Gravity}
\claimstatus{\StatusReduced}
\stagefield{Purpose}{Assemble the inter-defect force between two drain defects
from the matter (Noether) stress, earning the bilinear \(Q_1Q_2\) structure,
the reduced-3D \(r^{-2}\) and bulk-4D \(R^{-3}\) power laws, and the attractive
like-drain sign, consuming the Stage~001 geometry primitives.}
\stagefield{Verification}{SymPy audit:
\StageFile{scripts/ledger_stage002_matter_stress_force_assembly_sympy_audit.py}.
Mathematica audit:
\StageFile{mathematica/ledger_stage002_matter_stress_force_assembly_mathematica_audit.wl}.}

\paragraph{Status.}
\StatusReduced. Three verdict layers are carried from the earned source: the
leading matter-stress sign \texttt{FORCE\_ATTRACTIVE\_DERIVED} (target-blind);
the full sign \texttt{SIGN\_RESIDUAL\_QUANTUM\_VCONF\_MAXWELL\_PROFILE} (honest
residual); acceptance \texttt{PASS\_WITH\_NAMED\_RESIDUALS}. The bilinear
structure, the power laws, and the attractive sign are earned (form + sign); the
overall magnitude is a calibration knob; the full sign carries named,
profile/sim-deferred residuals.

\paragraph{Derivation.}
The medium obeys \(P=K\rho^5\), \(h=(5K/4)\rho^4\), with the barotropic identity
\(dP/d\rho=\rho\,(dh/d\rho)\). The matter (Noether) stress in Madelung form is
\[
  \Pi_{ij}=m\rho\,v_iv_j+\delta_{ij}P(\rho)+\sigma^{Q}_{ij},
  \qquad
  \sigma^{Q}_{ij}=\frac{\hbar^2}{4m}
  \Big[\frac{(\partial_i\rho)(\partial_j\rho)}{\rho}-\partial_i\partial_j\rho\Big],
\]
with the quantum-stress divergence contributing no net far-field force at this
order. A defect draining flux \(Q\) sets up, by Gauss's law, the inflow
\[
  v=-\frac{Q\,\hat n}{\Omega_{d-1}\,r^{\,d-1}},
\]
where \(\Omega_2=4\pi\) (\(d=3\)) and \(\Omega_3=2\pi^2\) (\(d=4\)) are the
Stage~001 solid angles. The Bernoulli cross-pressure of the two overlapping
inflows is \(\delta P=-m\rho\,(v_1\!\cdot\!v_2)\). Integrating the traction
\(\Pi_{ij}n_j\) over a control sphere around defect~2 and using the Stage~001
second moment \(\langle n_in_j\rangle=\delta_{ij}/d\) gives the convective factor
\(-(1+1/d)\) and the pressure factor \(+1/d\), whose sum is the total flux
\(-\,mN Q_2 v_1\) (the \(1/d\) pieces cancel). The stationary force
\(F_{12}=-\oint\Pi_{ij}n_j\,dS\), with defect~1's own inflow substituted, yields
the two lanes. \resultanchor{Reduced-3D force law.}
\[
  F_{12}=-\,\frac{m\,N_3\,Q_1Q_2}{4\pi\,r_{12}^2}\,\hat r_{12}.
\]
\resultanchor{Bulk-4D force law.}
\[
  F_{12}=-\,\frac{m\,\rho_4\,Q_1Q_2}{2\pi^2\,R_{12}^3}\,\hat R_{12}.
\]
\resultanchor{Attractive sign.} Since \(F\propto-Q_1Q_2\), like drains
(\(Q_1Q_2>0\)) attract and opposite signs repel; the sign is derived from the
Gauss-drain / Bernoulli / surface-integral chain, not assumed.

\paragraph{Physical use and named residuals.}
Target-blind, the inter-defect force is bilinear in the drain strengths, falls
as \(r^{-2}\) (reduced 3D) and \(R^{-3}\) (bulk 4D), and is attractive between
like drains. Honest scope: the overall normalization
(\(I_F/\Theta_Q/\text{branch-profile}\)) is a \emph{calibration} knob, not
derived; and the full sign carries the \texttt{QUANTUM\_STRESS\_FAR\_FIELD},
\texttt{VCONF\_BODY\_FORCE}, and \texttt{MAXWELL\_Z\_GAUGE\_JEXT} residuals,
which are profile/sim-deferred and first-class.
